\documentclass[pdftex,english,oribibl]{llncs}

%% Spracheinstellungen laden
\usepackage[english]{babel}

%% Schriftart in der Ausgabe/Eingabe
\usepackage[T1]{fontenc}
\usepackage{textcomp}
\usepackage[latin1]{inputenc}

%% Zitate
\usepackage[numbers]{natbib}
\bibliographystyle{abbrvnat}
%\bibliographystyle{dinat}
%\bibliographystyle{plainnat}
%\bibliographystyle{splncs}
%% Similar to option "sectionbib" but \refname instead of \bibname
\makeatletter
\renewcommand\bibsection{\section*{\refname\@mkboth{\MakeUppercase{\refname}}{\MakeUppercase{\refname}}}}
\makeatother

%% Index
%\usepackage{makeidx}
%\makeindex

%% PDF Einstellungen
% muss nach natbib geladen werden!
\usepackage{nameref}
\usepackage{varioref}
\usepackage[pdfusetitle,pdftex,colorlinks]{hyperref}
\hypersetup{pdfborder={0 0 0}}
\hypersetup{bookmarksdepth=3}
\hypersetup{bookmarksopen=true}
\hypersetup{bookmarksopenlevel=1}
\hypersetup{bookmarksnumbered=true}
\usepackage{color}
\hypersetup{colorlinks=false}
\usepackage{titling}

%\usepackage[section]{tocbibind}

\makeatletter
\gdef\@keywords{}
\def\keywords#1{\gdef\@keywords{#1}}
\gdef\@subtitle{}
\def\subtitle#1{\gdef\@subtitle{#1}}

%% modified from llncs
\renewenvironment{abstract}{%
  \list{}{\advance\topsep by0.35cm\relax\small%
          \leftmargin=1cm%
          \labelwidth=\z@%
          \listparindent=\z@%
          \itemindent\listparindent%
          \rightmargin\leftmargin}%
          \item[\hskip\labelsep\bfseries\abstractname]}{%
  \if!\@keywords!\else{\item[~]\item[\hskip\labelsep\bfseries\keywordname]\@keywords}\fi%
  \endlist}

\AtBeginDocument{%
  \if!\@subtitle!\else\hypersetup{pdfsubject={\@subtitle}}\fi
  \if!\@keywords!\else\hypersetup{pdfkeywords={\@keywords}}\fi
}
\makeatother

% llncs hyperref fix
\makeatletter
\providecommand*{\toclevel@author}{0}
\providecommand*{\toclevel@title}{0}
\makeatother

%% Grafiken
\usepackage[pdftex]{graphicx}
\DeclareGraphicsExtensions{.pdf,.jpg,.png}
\usepackage{subfigure}

%% Mathe
\usepackage{amsmath}
\usepackage{amssymb}

%% Listings
\usepackage{listings}
\lstset{escapechar=\%, frame=tb, basicstyle=\footnotesize}

%% Sonstiges
\newcommand{\TODO}[1]{\par\textcolor{red}{#1}\marginpar{\textcolor{red}{TODO}}}
\newcommand{\TODOX}[1]{\textcolor{red}{#1}\marginpar{\textcolor{red}{TODO}}}
\pagestyle{plain}

% Keine "Schusterjungen"
\clubpenalty = 10000
% Keine "Hurenkinder"
\widowpenalty = 10000 \displaywidowpenalty = 10000

%%%%%%%%%%%%%%%%%%%%%%%%%%%%%%%%%%%%%%%%%%%%%%%%%%%%%%%%%%%%%%%%%%%%%%%%%%%%%%%
%%% BEGIN DOCUMENT
%%%%%%%%%%%%%%%%%%%%%%%%%%%%%%%%%%%%%%%%%%%%%%%%%%%%%%%%%%%%%%%%%%%%%%%%%%%%%%%
\pretitle{%
  \begin{center}\LARGE
  \noindent\includegraphics[height=2.5cm]{figures/logo-se}\hfill{}\includegraphics[height=2.5cm]{figures/logo-uni}\\\vspace{0.5cm}
}
\title{Automatic debugging using comparable approaches}
% \subtitle{My (optional) Subtitle}
\author{Phillip, Daniel, Birgit Pohl (574353), Tim Sikatzki (573713)}
\institute{Humboldt University of Berlin\\Department of Computer Science\\12489 Berlin, Germany}
\posttitle{\end{center}}


\begin{document}

\maketitle

\begin{abstract}
  Context, Problematic, Solution

Astor is a automatic, multidisciplinary testing and repair library for Java. It contains three different modi, GenProg, Kali and MutRepair. Focussing on GenProg modus we found out that Astor's GenProg modus may not recognize enums, static variables that do not change their state.
To analyze Astor's effectivity we ran a test with Defects4J in order to find out, how many bugs is Astor able to fix. We conducted a bugfix on Astor and ran a test with the fixed version on Data4J as well. The results show .....
We also condcuted a test focussed on the test cases involving enums if they are repaired correctly.
\end{abstract}

\section{Introduction}
Definitions
Mentiones of relates work that is important in the introduction
Problem
- enums, static variables that do not change their state.
-  Dataset4J, a set of real-life peer-reviewed bugs for Java

Automatic Repair of Real Bugs: An Experience Report on
the Defects4J Dataset

Research Quetions
\subsection{RQ1}\label{rq1}

\section{Background}\label{sec:background}
Kommt alles das rein, was man braucht, um das Projekt zu verstehen. 

Our tool of choice is Astor a "publicly available automatic software repair tool"  for Java. Astor uses three different approaches, GenProg, Kali and MutRepair, which can be compared as well.

Fault Localisation
Repair Validation
Automatic Repair
GenProg
Kali
MutRepair



if statement nehmen und an richtige stelle einfügen, was er nicht macht
warum? varianz werden ausgegeben
tool löscht statements \cite{salam}

\section{Evaluation}\label{sec:evaluation}

\section{Solution and Evaluation Description}\label{sec:solutionEvaluationDescription}

\section{Discussion}\label{sec:discussion}

\section{Related Work}\label{sec:relatedWork}

\section{Conclusions}\label{sec:conclusions}

  These are my conclusions 

Dies ist ein ganz kurzer Beispieltext \cite{AspectJ2007}

\section{References}\label{sec:references}
\bibliography{Report-Astor}







\end{document}
